\scanpage{109}
% L1--8; oblique cancellation strokes through the entire first corollary.
\deleted{\uline{Corollaire} 1 Supposons que dans $X$ et $Y$, \deleted{tt pt} \unsure{admettant|admettent} une \deleted{\unsure{application fermée est de même}} \inserted{base d'ouverts \inserted{$X_i$ resp. $Y_j$} rétrocompacts quasi-compacts.} Alors l'hyp préliminaire de la prop. équivaut aussi à dire que pour tt ouvert \deleted{\unsure{qu. cpt}} $Y_j$, $f^{-1}(Y_j)$ est rétrocompact dans $X$, ou encore que pour tt \deleted{ouvert} ouvert $Y'$ \deleted{\unsure{qu. cpt dans}} \inserted{rétrocompact \deleted{\unsure{(ou rétrocompact)}}} $Y$, $X'=f^{-1}(Y')$ est \deleted{qu. cpt} rétrocompact dans $X$, ou encore que les $X_i\cap f^{-1}(Y_j)$ sont qu. cpts.}

% L9--13
[ \uline{Prop.} Soit $X$ un espace topologique, ayant une base d'ouverts $U_i$ qui sont quasi-compacts. Alors pour tte \deleted{partie $E$ de $X$, et \unsure{ouvert|toute}} application continue $f:X'\to X$, il revient au $=$ de dire que $f$ est qu. cpt, ou que les $f^{-1}(U_i)$ sont qu. cpts.

% L14--16
Donc si $X'\subset X$ sous-espace top., $X'$ est rétrocompact dans $X$ ssi les $X'\cap U_i$ sont qu. cpts.

% L17--19
\uline{Cor.} \deleted{Soit Supposons de \uline{plus}} \inserted{Pour que} les $U_i$ rétrocompacts dans $X$, \deleted{Alors} \inserted{il f. et s. que} \uuline{tte} partie quasi-compacte ouverte de $X$ soit rétrocompacte dans $X$.

% L20--25
\uline{Cor} S'il en est ainsi, pour une \deleted{\unsure{m}} application continue $f:X'\to X$ il revient au $=$ de dire \deleted{\unsure{que les}} (a) pour tt \inserted{ouvert} $U\subset X$ qui est rétrocompact, $f^{-1}(U)$ est rétrocompact, (b) Pour tt $U_i$, $f^{-1}(U_i)$ est rétrocompact dans $X'$, (c) [Si $X'$ a une base d'ouverts \deleted{\unsure{quasi rétrocompacts|quasicompacts}} $U'_j$] les $f^{-1}(U_i)\cap U'_j$ sont \deleted{rétro} quasi cpts.]

\margin{\deleted{\uline{Cor} Si $f$ \inserted{$:X'\to X$} est qu. cpt, et \unsure{dans $X'$ ttes les parties}}}
% slanted

% L22--27, small proof at lower left, continuing across the foot.
\margin{$b\Rightarrow a$, En effet, \inserted{$U\subset X$ ouv. rétrocomp., si $U'\subset X'$ ouv. qu. cpt} $f^{-1}(U)\cap V'$ \deleted{$=f^{-1}(U)\cap \unsure{f^{-1}}\cap$} $\subset f^{-1}(f(\unsure{U'|V'}))$, et $f(V')$ étant qu. cpt, est contenu dans une réunion \uline{finie} de $U_i$, $i\in J$, donc $f^{-1}(U)\cap V'=f^{-1}(\bigcup_{i\in J} U\cap U_i)\cap V'$ Or chaque $U\cap U_i$ est réunion \uline{finie} \unsure{d'ens|d'une|d'un} $\unsure{U_j|U_l}$}
% archive sheet number: 108
